\documentclass[11pt]{article}
\usepackage[margin=1in]{geometry}
\usepackage{setspace}
\usepackage{parskip}

\begin{document}
\newpage

\begin{flushleft}
\textbf{Letter 1: Recommendation for Luke Brown}

\vspace{12pt}

To Whom It May Concern,

\vspace{12pt}

It is my pleasure to recommend Mr. Luke Brown, who served as the Signal Processing Lead and Team Leader on our capstone project. Luke demonstrated strong technical skills and effective leadership throughout the semester.

Luke’s Diligence was \textbf{Very Good}. He consistently completed his tasks ahead of schedule and with reliable, positive results. A notable example was his work on optimizing the GPS signal acquisition algorithm, where his proactive debugging reduced processing time by 30\%. His dedication ensured our transition from MATLAB to Python progressed smoothly, allowing the team to leverage more efficient programming practices and enhance overall system performance.

In terms of Cooperativeness, Luke was \textbf{Excellent}. He was always present at meetings, guided discussions towards consensus, and frequently offered assistance to teammates. For instance, during a debate about resource allocation for FFT-based correlation, Luke mediated effectively, balancing the needs of both the GUI and antenna teams. His ability to foster a collaborative environment ensured that all team members felt valued and heard, which significantly contributed to the project's harmonious progression.

Luke’s Resourcefulness was \textbf{Very Good}. When faced with connectivity issues with the RTL-SDR v4, he identified and implemented a solution by switching to the RTL-SDR v3, minimizing delays. Additionally, his discovery of open-source libraries optimized our data throughput without exceeding the project budget. His knack for finding cost-effective and efficient solutions was invaluable, especially when navigating the complexities of real-time signal processing.

Luke’s Communicativeness was \textbf{Excellent}. He provided clear and concise documentation for the signal processing workflow, making it accessible for all team members. His presentations on phase-locked loop debugging were particularly insightful and facilitated the identification of critical issues. Furthermore, Luke maintained open lines of communication, ensuring that any potential misunderstandings were promptly addressed and that the team remained aligned with project goals.

Luke’s overall Productivity was approximately \textbf{20\%}. While slightly below an equal distribution, his contributions in leadership and technical problem-solving were integral to the project’s success. His ability to manage both the technical and interpersonal aspects of the team ensured that our milestones were consistently met and that the project advanced smoothly.

I strongly recommend Luke for positions requiring technical expertise and team leadership. His combination of technical acumen, proactive problem-solving, and exceptional leadership qualities make him a valuable asset to any engineering team.

\vspace{12pt}

Sincerely, \\
J. C. Vaught

\end{flushleft}

\newpage

\begin{flushleft}
\textbf{Letter 2: Recommendation for Steven Schweninger}

\vspace{12pt}

To Whom It May Concern,

\vspace{12pt}

I am pleased to recommend Mr. Steven Schweninger, who served as the GUI Developer and Integration Specialist for our capstone project. Steven’s work was vital to ensuring the system’s usability and functionality, playing a crucial role in the seamless integration of various subsystems.

Steven’s Diligence was \textbf{Very Good}. He consistently completed his tasks ahead of schedule. For example, he developed a graphics driver for the NewHaven OLED screen, resolving compatibility issues with the SSD1333 chip and enabling real-time GPS visualization. His ability to anticipate project needs and deliver high-quality work promptly ensured that the user interface was both functional and aesthetically pleasing, meeting all project requirements effectively.

Steven’s Cooperativeness was \textbf{Very Good}. He participated in every meeting, contributed to group decisions, and assisted teammates when needed. For instance, he helped the antenna team debug a formatting issue with GPS coordinates, ensuring seamless integration with the GUI. His willingness to collaborate and support other team members fostered a cooperative team environment, enhancing overall project cohesion and efficiency.

Steven’s Resourcefulness was \textbf{Very Good}. He identified and integrated Folium for map rendering and Selenium for static image generation, overcoming the lack of built-in internet-free map capabilities. These innovative solutions ensured the GUI met sponsor requirements for real-time map displays. Additionally, his ability to adapt and implement new technologies significantly enhanced the functionality and user experience of our system.

Steven’s Communicativeness was \textbf{Very Good}. His updates on GUI parameters were timely and detailed, allowing other team members to understand how to integrate their subsystems effectively. He maintained clear and open communication channels, ensuring that any changes or updates to the GUI were promptly shared and understood by all relevant parties. This clarity in communication was essential for the smooth integration and operation of the overall system.

Steven contributed approximately \textbf{25\%} of the team’s overall progress. His focus on usability and attention to detail significantly enhanced the user experience, making the system more intuitive and user-friendly. His dedication to creating a robust and reliable GUI was a key factor in the project's success, ensuring that end-users could interact with the system efficiently and effectively.

I highly recommend Steven for roles that require technical expertise in GUI design and system integration. His ability to blend technical skills with a collaborative spirit makes him an excellent candidate for any team-oriented engineering position.

\vspace{12pt}

Sincerely, \\
J. C. Vaught

\end{flushleft}

\newpage

\begin{flushleft}
\textbf{Letter 3: Recommendation for Parker Huggins}

\vspace{12pt}

To Whom It May Concern,

\vspace{12pt}

I am pleased to recommend Mr. Parker Huggins, who served as the RF/SDR Lead on our capstone project. Parker’s deep technical expertise and problem-solving skills were crucial to the team’s progress and the successful implementation of our SDR subsystem.

Parker’s Diligence was \textbf{Excellent}. He consistently delivered results ahead of schedule, such as configuring the RTL-SDR v3 with a low-noise amplifier and bias tee to enhance signal quality. His proactive efforts ensured the SDR subsystem met performance benchmarks during testing, allowing the team to maintain our project timeline without compromising on quality.

Parker’s Cooperativeness was \textbf{Very Good}. He attended nearly all meetings, contributed to team discussions, and assisted others with subsystem integration. When the antenna team encountered impedance matching issues, Parker provided valuable guidance and code snippets to resolve the problem efficiently. His collaborative approach ensured that technical challenges were addressed promptly, facilitating smooth subsystem integration and overall project harmony.

Parker’s Resourcefulness was \textbf{Excellent}. He navigated complex technical documentation and adapted GNSS-SDRLIB for real-time processing, significantly enhancing our receiver’s capabilities. His suggestion to implement a Butterworth low-pass filter improved signal integrity and reduced noise, demonstrating his ability to apply theoretical knowledge to practical challenges. Parker’s innovative solutions were instrumental in overcoming technical hurdles and optimizing system performance.

Parker’s Communicativeness was \textbf{Very Good}. While his communications were sometimes brief, they were always timely and clear. During presentations, he concisely explained the intricacies of phase-locked loops, ensuring the team presented a cohesive understanding of the SDR subsystem. His ability to convey complex technical concepts in an understandable manner was essential for team alignment and effective project execution.

Parker’s Productivity was approximately \textbf{30\%}, reflecting his pivotal role in the project. His expertise and reliability were key to achieving our goals, particularly in the areas of RF system optimization and SDR integration. Parker’s substantial contributions not only advanced the technical aspects of the project but also supported the overall team dynamic and project management.

I highly recommend Parker for roles requiring advanced technical skills and problem-solving in RF systems. His combination of technical proficiency, resourcefulness, and collaborative spirit make him an outstanding candidate for any engineering team focused on cutting-edge technology and system integration.

\vspace{12pt}

Sincerely, \\
J. C. Vaught

\end{flushleft}

\newpage

\begin{flushleft}
\textbf{Letter 4: Cover Letter for Myself (J.C. Vaught)}

\vspace{12pt}

To Whom It May Concern,

\vspace{12pt}

My name is J.C. Vaught, and I served as the Antenna Lead and Project Scribe for our SDR-based GNSS receiver project. My role focused on designing and testing antennas to optimize signal acquisition for real-time operations, as well as documenting the project’s progress and outcomes.

My Diligence was \textbf{Satisfactory}. I completed all assigned tasks on time, including generating 3D radiation patterns and gain measurements for the L1 GPS band. These results guided the team in selecting a 28 dB gain antenna that balanced performance with cost-effectiveness. While I consistently met deadlines, there were instances where additional feedback from teammates was beneficial to refine my deliverables further.

My Cooperativeness was \textbf{Very Good}. I actively participated in team discussions, providing data-driven input for trade-offs. For example, I collaborated with the SDR lead to validate impedance matching, ensuring subsystem compatibility. My willingness to engage with team members and contribute to collective decision-making processes fostered a collaborative and supportive team environment, enhancing our project’s overall effectiveness.

My Resourcefulness was \textbf{Satisfactory}. While I initially required guidance for certain test setups, I independently developed a rotating platform for radiation pattern testing and adapted methods to measure power draw using limited equipment. This ability to innovate and adapt under resource constraints demonstrated my capacity to contribute effectively, even when faced with challenges or limited resources.

My Communicativeness was \textbf{Very Good}. I consistently shared updates on antenna performance metrics and clarified requirements with other subsystem leads, ensuring smooth integration across the team. For example, I provided detailed reports on antenna testing results, which informed the SDR and signal processing teams’ subsequent adjustments. My clear and timely communication helped prevent misunderstandings and kept the project on track.

My Productivity was approximately \textbf{25\%}, reflecting my contributions to subsystem design, testing, and documentation. This included about 20\% from written reports and 5\% from subsystem-specific work. By offering reliable antenna design insights and thorough documentation, I helped ensure that the team could present a robust, well-tested system. Additionally, my role as Project Scribe ensured that all project activities were well-documented, facilitating transparency and accountability within the team.

This project enhanced my skills in collaborative problem-solving and experimental design, which I am eager to apply to future challenges. I am confident that the experiences gained from this capstone project have prepared me to contribute meaningfully to your team, bringing a strong foundation in antenna design, technical documentation, and teamwork.

\vspace{12pt}

Sincerely, \\
J.C. Vaught

\end{flushleft}

\end{document}